% =====================================================================
%  2bat-ud1-13.tex  —  SESSIÓ 13: Quan es pot multiplicar? Compatibilitat
%  (sense preàmbul: s'inclou des de main.tex)
%  Criteris: C3 (producte) i C4 (interpretar). Ampliació conceptual.
% =====================================================================
\horatitol{Sessió 13 — Quan es pot multiplicar? Compatibilitat i interpretació}%
{Decidir si un producte de matrices té sentit, anticipar-ne la dimensió i interpretar què representa el resultat.}

\subsection{Idea clau}

Saber \emph{calcular} un producte de matrices no és tot: cal saber \textbf{quan es
pot fer} i \textbf{què significa} el resultat. De fet, decidir si un producte té
sentit i interpretar-lo és tan important com fer els números. Aquesta sessió posa
el focus en aquestes dues coses.

\begin{idea}
\textbf{Compatibilitat, dimensió i interpretació}
\begin{itemize}[leftmargin=1.4em, itemsep=2pt]
  \item \textbf{Compatibilitat:} per multiplicar $A\cdot B$, el nombre de
        \textbf{columnes} de $A$ ha de coincidir amb el nombre de \textbf{files}
        de $B$. Si no, el producte \emph{no es pot fer}.
  \item \textbf{Dimensió del resultat:} si $A$ és $m\times n$ i $B$ és $n\times p$,
        el resultat $A\cdot B$ és $m\times p$ (files de la primera, columnes de la
        segona).
  \item \textbf{Regla visual:} $\;(m\times \underline{n})\cdot(\underline{n}\times p)
        \rightarrow m\times p$. Les dues $n$ del mig han de coincidir; les de fora
        donen la dimensió del resultat.
  \item \textbf{Interpretació:} les files del resultat hereten el \emph{significat}
        de les files de la primera matriu; les columnes, el de les columnes de la
        segona.
\end{itemize}
\end{idea}

\subsection{Exemples}

\begin{exemple}[quins productes es poden fer]
Mirem només les dimensions:
\begin{itemize}[leftmargin=1.4em, itemsep=2pt]
  \item $(2\times 3)\cdot(3\times 4)$: les dues $3$ coincideixen $\rightarrow$
        \textbf{es pot fer}, resultat $2\times 4$.
  \item $(2\times 3)\cdot(2\times 3)$: $3\neq 2$ $\rightarrow$ \textbf{no es pot fer}.
  \item $(3\times 1)\cdot(1\times 2)$: les dues $1$ coincideixen $\rightarrow$
        \textbf{es pot fer}, resultat $3\times 2$.
\end{itemize}
\end{exemple}

\begin{exemple}[què representa el resultat]
Una matriu $U$ recull, per a cada usuari (files), quantes sessions de cada servei
(columnes) necessita: és $3\times 2$ (3 usuaris, 2 serveis). Una matriu $C$ recull,
per a cada servei (files), el cost i les hores (columnes): és $2\times 2$.
El producte $U\cdot C$ és compatible ($2=2$) i dona una matriu $3\times 2$: per a
cada \textbf{usuari} (files, heretades de $U$), el \textbf{cost} i les \textbf{hores}
totals (columnes, heretades de $C$).
\end{exemple}

\subsection{Exercicis resolts}

\begin{exresolt}[decidir, anticipar i interpretar]
Una matriu $A$ de dimensió $4\times 3$ recull, per a quatre centres (files), les
places de tres serveis (columnes). Una matriu $B$ de dimensió $3\times 2$ recull,
per a cada servei (files), el cost i el personal necessari (columnes).
\begin{enumerate}[label=\alph*), leftmargin=1.6em, topsep=2pt]
  \item Es pot calcular $A\cdot B$? Quina dimensió tindria?
  \item Què representarien les files i les columnes del resultat?
  \item Es pot calcular $B\cdot A$?
\end{enumerate}
\solucio
a) $A$ és $4\times 3$ i $B$ és $3\times 2$: les columnes de $A$ ($3$) coincideixen
amb les files de $B$ ($3$). Sí, es pot fer, i el resultat és $4\times 2$.

b) Les files (heretades de $A$) són els \textbf{centres}; les columnes (heretades
de $B$) són el \textbf{cost} i el \textbf{personal} totals de cada centre.

c) $B$ és $3\times 2$ i $A$ és $4\times 3$: les columnes de $B$ ($2$) no
coincideixen amb les files de $A$ ($4$). $B\cdot A$ \emph{no} es pot fer.
\resposta{$A\cdot B$ sí (dimensió $4\times 2$): files $=$ centres, columnes $=$ cost i personal totals. $B\cdot A$ no es pot fer.}
\end{exresolt}

\begin{exresolt}[l'ordre canvia el significat]
Amb $U=\begin{pmatrix} 2 & 1 \\ 1 & 3 \end{pmatrix}$ (usuaris $\times$ tipus de
sessió) i $D=\begin{pmatrix} 2 & 10 \\ 1 & 8 \end{pmatrix}$ (tipus de sessió
$\times$ hores i cost), calcula $U\cdot D$ i explica per què $D\cdot U$, tot i que
es pot calcular, \emph{no} tindria el mateix significat.
\solucio
\[
U\cdot D=\begin{pmatrix} 2\per 2+1\per 1 & 2\per 10+1\per 8 \\ 1\per 2+3\per 1 & 1\per 10+3\per 8 \end{pmatrix}
=\begin{pmatrix} 5 & 28 \\ 5 & 34 \end{pmatrix}.
\]
Aquí les files són els usuaris i les columnes, les hores i el cost: és la lectura
que té sentit. En $D\cdot U$ es creuarien tipus de sessió amb usuaris d'una manera
que no respon a cap pregunta real; encara que els números surtin, no representen
res interpretable. Per això l'ordre dels factors importa també en el \emph{significat}.
\resposta{$U\cdot D=\begin{pmatrix} 5 & 28 \\ 5 & 34 \end{pmatrix}$ (usuaris $\times$ hores i cost). $D\cdot U$ no té lectura social.}
\end{exresolt}

\subsection{Exercicis per fer}

\begin{exercicis}
\begin{exlist}
  \item Digues, \emph{només mirant les dimensions}, quins d'aquests productes es
        poden fer i quina dimensió tindria el resultat:
        \begin{enumerate}[label=\alph*), leftmargin=1.6em, topsep=2pt]
          \item $(3\times 2)\cdot(2\times 5)$
          \item $(2\times 4)\cdot(2\times 4)$
          \item $(1\times 3)\cdot(3\times 1)$
          \item $(4\times 1)\cdot(1\times 3)$
        \end{enumerate}
  \item Una matriu d'usuaris per serveis és $5\times 3$ i una de serveis per cost
        és $3\times 1$. Es pot fer el producte? Quina dimensió tindria i què
        representaria?
  \item Per a cadascun, digues si $A\cdot B$ es pot fer, si $B\cdot A$ es pot fer,
        o totes dues coses:
        \begin{enumerate}[label=\alph*), leftmargin=1.6em, topsep=2pt]
          \item $A$ és $2\times 3$, $B$ és $3\times 2$
          \item $A$ és $2\times 2$, $B$ és $2\times 2$
          \item $A$ és $3\times 2$, $B$ és $2\times 4$
        \end{enumerate}
  \item Una matriu $T$ de tallers per casal és $2\times 3$ (2 tallers, 3 casals) i
        una matriu $H$ d'hores és $3\times 1$. Calcula $T\cdot H$ amb
        $T=\begin{pmatrix} 2 & 1 & 3 \\ 1 & 4 & 0 \end{pmatrix}$ i
        $H=\begin{pmatrix} 2 \\ 1 \\ 3 \end{pmatrix}$, i digues què representa cada
        número del resultat.
  \item Una companya vol multiplicar una matriu de places ($3\times 4$) per una de
        costos ($3\times 1$) per saber el cost de cada centre. Per què no funciona?
        Com hauria de ser la matriu de costos perquè el producte tingués sentit?
  \item \textbf{(Raonament)} Tens una matriu $A$ de dimensió $2\times 3$. Quina
        dimensió ha de tenir una matriu $B$ perquè $A\cdot B$ es pugui fer i el
        resultat sigui una matriu $2\times 2$? Justifica-ho amb la regla visual.
\end{exlist}
\end{exercicis}

% ---------------------------------------------------------------------
\fullsolucions{Sessió 13}
\begin{solucions}
\begin{sollist}
  \item (a) es pot fer, $3\times 5$ \quad (b) no es pot fer ($4\neq 2$) \quad
        (c) es pot fer, $1\times 1$ \quad (d) es pot fer, $4\times 3$.
  \item Sí: $5\times 3$ per $3\times 1$ dona $5\times 1$. Representaria, per a cada
        usuari (files), el cost total (única columna).
  \item (a) totes dues ($A\cdot B$ és $2\times 2$, $B\cdot A$ és $3\times 3$) \quad
        (b) totes dues (les dues $2\times 2$) \quad (c) només $A\cdot B$ ($3\times 4$);
        $B\cdot A$ no ($4\neq 3$).
  \item $T\cdot H=\begin{pmatrix} 2\per 2+1\per 1+3\per 3 \\ 1\per 2+4\per 1+0\per 3 \end{pmatrix}
        =\begin{pmatrix} 14 \\ 6 \end{pmatrix}$. Cada número són les hores totals
        d'un taller sumant els tres casals.
  \item No funciona perquè $(3\times 4)\cdot(3\times 1)$ té $4\neq 3$: no és
        compatible. Perquè tingui sentit, la matriu de costos ha de ser
        $4\times 1$ (un cost per cada servei), de manera que
        $(3\times 4)\cdot(4\times 1)$ doni un cost per centre ($3\times 1$).
  \item $B$ ha de ser $3\times 2$: la regla visual és
        $(2\times \underline{3})\cdot(\underline{3}\times 2)\rightarrow 2\times 2$.
        Les columnes de $A$ ($3$) i les files de $B$ han de coincidir, i les files
        de $A$ ($2$) amb les columnes de $B$ ($2$) donen el resultat.
\end{sollist}
\end{solucions}
